\documentclass[10pt,a4paper]{amsart}
\usepackage[margin=19mm]{geometry}
\usepackage{amsmath,amssymb,mathtools}
\usepackage[scr=rsfs,cal=euler]{mathalfa}
\usepackage{xfrac}
\usepackage[unicode,hidelinks,bookmarks=false]{hyperref}
\newcommand{\ie}{i.e.\ }
\newcommand{\cf}{cf.\ }
\newcommand{\resp}{respectively\ }
\newcommand{\Subsection}{\subsection}
\providecommand{\nobreakdash}{\nobreak\hbox{-}}
\setlength{\emergencystretch}{2em}
\multlinegap0pt
\usepackage{amssymb}
\usepackage[shortlabels]{enumitem}
\newtheorem{lemma}{Lemma}
\usepackage{cleveref}
\crefname{lemma}{Lemma}{Lemmas}
\newcommand{\pres}[2]{\left\langle #1\ \middle|\ #2\right\rangle} % group presentation
\title{Hypercubical manifolds in homotopy type theory}
\author{Samuel Mimram, \'Emile Oleon}
\date{Source: arXiv:2506.19402v2 (CC BY 4.0)}
\begin{document}
\maketitle

\noindent\emph{Source and licence.} Hypercubical manifolds in homotopy type theory. Version arXiv:2506.19402v2 (CC BY 4.0), distributed by arXiv under the Creative Commons Attribution 4.0 International licence, http://creativecommons.org/licenses/by/4.0/, as recorded in the arXiv licence metadata for this exact version and shown on the abstract page https://arxiv.org/abs/2506.19402v2. Full licence notice: \texttt{licenses/arxiv-2506.19402-CC-BY-4.0.txt}.\par\smallskip

\noindent\textit{Editorial scope.} Counted unit: the article's lemma q-pres (source0.tex lines 103--110). The article's complete original proof of it is delivered with it (source0.tex lines 111--151, one \texttt{proof} environment of 41 lines). No mathematics is written, repaired or re-derived: the statement and the proof are the article's own lines, in the article's own notation, and no retained line is substituted or re-flowed.  No further source-local context is needed: every cross-reference of every retained line resolves inside the two delivered spans.  Nothing was omitted from any retained span, so the package carries no omission record.  No retained line contains a citation, so the wrapper carries no bibliography and no bibliography entry.  No retained line contains a figure environment, an image-inclusion command or drawing code, so no image is delivered and the package declares no asset dependency.  Editorial formatting only, in addition: an AMS \texttt{amsart} preamble that restates the article's own declarations and macros used by the retained lines, namely the environments \texttt{lemma} on separate counters, together with the standard packages those lines need.  The retained environments print in this document as Lemma 1 in order of appearance, whereas the article prints the same items with the numbering of its own full document.  The complete native source of the article as distributed in the arXiv e-print for this exact version is bundled unchanged as \texttt{sources/180413/source0.tex}.
\begin{lemma}
  \label[lemma]{quaternion-pres}
  The quaternion group admits the presentation
  \begin{equation}
    \label{q-pres}
    Q=\pres{i,j}{i=jij,j=iji}
  \end{equation}
\end{lemma}

\begin{proof}
  We first show that the group admits the presentation
  \begin{equation}
    \label{q-pres2}
    Q=\pres{i,j}{i^4=1,i^2=j^2,iji=j}
  \end{equation}
  We have the following equivalent presentations
  \begin{align*}
    Q&=\pres{e,i,j,k}{i^2=e,j^2=e,k^2=e,ijk=e,e^2=1}
    \\
    Q&=\pres{i,j,k}{j^2=i^2,k^2=i^2,ijk=i^2,i^4=1}&&\text{using $e=i^2$}
  % Above, $ij=k$ is derivable by
  % \begin{align*}
    % ijk&=i^2\\
    % ijk^4&=i^2k^3&&\text{by multiplication by $k^3$}\\
    % ij&=k&&\text{by $k^4=i^4=1$ and $i^2k^2=i^4=1$.}
    \\
    \intertext{
      In the second presentation above, $ij=k$ is derivable because we have $ijk=i^2$, multiplying by~$k^3$ yields $ijk^4=i^2k^3$, and therefore $ij=k$ since $k^4=i^4=1$ and $i^2k^2=i^4=1$. We thus have the following presentations:
      }
  % \begin{align*}
    Q&=\pres{i,j,k}{j^2=i^2,k^2=i^2,ijk=i^2,i^4=1,ij=k}%&&\text{$ij=k$ derivable}
    \\
    Q&=\pres{i,j}{j^2=i^2,(ij)^2=i^2,ijij=i^2,i^4=1}&&\text{using $k=ij$}
    \\
    \intertext{
      The relation $(ij)^2=i^2$ is now redundant, so that we can remove it, and then, from the relation $ijij=i^2=j^2$, we can derive the equivalent relation $iji=j$ by dividing by $j$, and we obtain \cref{q-pres2}:
    }
    Q&=\pres{i,j}{i^4=1,i^2=j^2,iji=j}&&\text{\qquad\qquad\qquad\qquad\qquad\qquad\qquad}
  \end{align*}

  We can now show that the presentation \cref{q-pres2} is equivalent to \cref{q-pres}. From \cref{q-pres2}, we can derive the missing relation $i=jij$
  % since $jij$ is equal to $iji^2j$ by $j=iji$, and thus to $ij^4$ by $i^2=j^2$ and thus to $i$ by $j^4=i^4=1$.
  by
  \begin{align*}
    jij&=iji^2j&&\text{by $j=iji$}\\
    &=ij^4&&\text{by $i^2=j^2$}\\
    &=i&&\text{by $j^4=i^4=1$}
  \end{align*}
  Conversely, from \cref{q-pres} we can derive the relations of \cref{q-pres2} since $i^2=jiji=j^2$, and $j=iji=jij^2i=ji^4$ which implies $1=i^4$ by division by~$j$.
\end{proof}

\end{document}
